\documentclass{ctexart}
\usepackage{geometry}
\usepackage{fancyhdr}
\usepackage{graphicx}
\usepackage{booktabs}
\usepackage{amsmath}
\usepackage{tikz}
\usepackage{array}
\xeCJKsetup{CJKmath=true} 
\usepackage{zhnumber} % change section number to chinese
\renewcommand\thesection{\zhnum{section}}
\renewcommand \thesubsection {\arabic{subsection}}
\CTEXsetup[format={\Large\bfseries}]{section}

\geometry{
    a4paper,
    left=3.18cm,
    right=3.18cm,
    top=3.04cm,
    bottom=3.04cm
}

\pagestyle{fancy}
\fancyhf{}
\renewcommand{\headrulewidth}{0.7pt} % 设置页眉横线粗细
\fancyhead[L]{\kaishu\large 大学物理实验报告} % 在左侧设置页眉文字
\fancyhead[R]{\kaishu\large 哈尔滨工业大学(深圳) } % 在右侧设置页眉文字
\fancyfoot[R]{\thepage} % 将页数放在右下角


\setlength\headwidth{\textwidth}

\begin{document}

\noindent
\begin{center}
\textbf{
\begin{tabular}{p{2.4cm}p{2.4cm}p{4cm}p{4cm}}
    班级 \hrulefill & 学号 \hrulefill & 姓名 \hrulefill & 教师签字 \hrulefill \\
\end{tabular}
\begin{tabular}{p{6cm}p{3.6cm}p{3.6cm}}
    实验日期 \hrulefill & 预习成绩 \hrulefill & 总成绩 \hrulefill
\end{tabular}
{\noindent}	 \rule[-10pt]{\textwidth}{0.7pt}
}\end{center}

\begin{center}
    \Large \textbf{实验内容 \underline{磁耦合共振式无线电力传输实验}}
\end{center}


\section{实验预习指导}
略
\newpage

\section{原始数据记录}
略
\newpage

\section{数据处理}
\subsection{研究震荡频率对电力传输效率的影响}
绘制幅度-频率曲线，总结曲线规律。

\subsection{研究无线电力传输距离对传输效果的影响}
绘制灯泡电压-距离曲线，总结曲线规律。

\subsection{自制无线电力传输系统}
总结实际传输效果，分析误差产生的原因。

\section{讨论题}
\subsection{为什么当震荡频率和LC电路的频率一样时，发射线圈能在周围产生大的交变磁场？}

\subsection{你认为提高磁耦合谐振式无线电力传输系统能量传输效率的方式有哪些？}

\end{document}
